% GENERATED FILE. Edit canonical lessons, then run npm run generate:books.
% canonical-source-hash: fnv1a64:e0abfe345fd50084
% canonical-lessons: AR-C97-ashara, AR-C97-thalathun, AR-C97-arbaun, AR-C97-khamsun, AR-C97-alf

\chapter{Ten, Thirty, Forty, Fifty, Thousand}
\label{ch:ar-ten-thirty-forty-fifty-thousand}
\hlchaptermodality{97}

\begin{chapteropening}
\textbf{By the end of this chapter:} \emph{I can say ten, thirty, forty, fifty and a thousand.}

Complete the last lesson of chapter 97: I can say ten, thirty, forty, fifty and a thousand.
\end{chapteropening}

\section[\texorpdfstring{\ar{عشرة} (ʿashara)}{عشرة (ʿashara)}]{\ar{عشرة} (ʿashara) — ten}

\label{lesson:AR-C97-ashara}



\begin{quote}
Before the new one: say the Arabic for eight, then the Arabic for nine.
\end{quote}

\subsection*{You'll want to know: \ar{عشرة}}
\textbf{\ar{عشرة}} — \emph{ʿashara} — \textquotedblleft{}ten\textquotedblright{}.

\begin{cousinweb}
\textbf{\ar{عشرة}} (\emph{ʿashara}) is built on the native Arabic root \emph{ʿ-sh-r}, \textquotedblleft{}ten\textquotedblright{} (Lane). Hebrew \emph{ʿeser}, \textquotedblleft{}ten\textquotedblright{}, is its cousin.
\end{cousinweb}

\begin{grammarlens}[title={What you've built}]
One more word for this chapter.
\end{grammarlens}

\subsection*{Your turn}
\begin{itemize}
\raggedright
  \item \emph{Say it:} \emph{ʿashara}
  \item \emph{Say it:} \emph{ʿashara}, once more
  \item \emph{Say it:} say \emph{ʿashara}
  \item \emph{Recall:} say the Arabic for eight, then the Arabic for nine, then say \emph{ʿashara} again
\end{itemize}

\begin{tcolorbox}[breakable,colback=teal!4,colframe=teal!35!black,title={Before you move on}]
What does \ar{عشرة} mean? (\textquotedblleft{}Ten\textquotedblright{}.) Say it once more.
\end{tcolorbox}
\section[\texorpdfstring{\ar{ثلاثون} (thalāthūn)}{ثلاثون (thalāthūn)}]{\ar{ثلاثون} (thalāthūn) — thirty}

\label{lesson:AR-C97-thalathun}



\begin{quote}
Before the new one: say the Arabic for nine, then the Arabic for ten.
\end{quote}

\subsection*{You'll want to know: \ar{ثلاثون}}
\textbf{\ar{ثلاثون}} — \emph{thalāthūn} — \textquotedblleft{}thirty\textquotedblright{}.

\begin{cousinweb}
\textbf{\ar{ثلاثون}} (\emph{thalāthūn}) is built on the native Arabic root \emph{th-l-th}, \textquotedblleft{}three\textquotedblright{} (Lane). \emph{Thalāthūn} is \emph{thalāth}, \textquotedblleft{}three\textquotedblright{}, with the tens ending \emph{-ūn}.
\end{cousinweb}

\begin{grammarlens}[title={What you've built}]
One more word for this chapter.
\end{grammarlens}

\subsection*{Your turn}
\begin{itemize}
\raggedright
  \item \emph{Say it:} \emph{thalāthūn}
  \item \emph{Say it:} \emph{thalāthūn}, once more
  \item \emph{Say it:} say \emph{thalāthūn}
  \item \emph{Recall:} say the Arabic for nine, then the Arabic for ten, then say \emph{thalāthūn} again
\end{itemize}

\begin{tcolorbox}[breakable,colback=teal!4,colframe=teal!35!black,title={Before you move on}]
What does \ar{ثلاثون} mean? (\textquotedblleft{}Thirty\textquotedblright{}.) Say it once more.
\end{tcolorbox}
\section[\texorpdfstring{\ar{أربعون} (arbaʿūn)}{أربعون (arbaʿūn)}]{\ar{أربعون} (arbaʿūn) — forty}

\label{lesson:AR-C97-arbaun}



\begin{quote}
Before the new one: say the Arabic for ten, then the Arabic for thirty.
\end{quote}

\subsection*{You'll want to know: \ar{أربعون}}
\textbf{\ar{أربعون}} — \emph{arbaʿūn} — \textquotedblleft{}forty\textquotedblright{}.

\begin{cousinweb}
\textbf{\ar{أربعون}} (\emph{arbaʿūn}) is built on the native Arabic root \emph{r-b-ʿ}, \textquotedblleft{}four\textquotedblright{} (Lane). \emph{Arbaʿūn} is \emph{arbaʿ}, \textquotedblleft{}four\textquotedblright{}, with the tens ending \emph{-ūn}.
\end{cousinweb}

\begin{grammarlens}[title={What you've built}]
One more word for this chapter.
\end{grammarlens}

\subsection*{Your turn}
\begin{itemize}
\raggedright
  \item \emph{Say it:} \emph{arbaʿūn}
  \item \emph{Say it:} \emph{arbaʿūn}, once more
  \item \emph{Say it:} say \emph{arbaʿūn}
  \item \emph{Recall:} say the Arabic for ten, then the Arabic for thirty, then say \emph{arbaʿūn} again
\end{itemize}

\begin{tcolorbox}[breakable,colback=teal!4,colframe=teal!35!black,title={Before you move on}]
What does \ar{أربعون} mean? (\textquotedblleft{}Forty\textquotedblright{}.) Say it once more.
\end{tcolorbox}
\section[\texorpdfstring{\ar{خمسون} (khamsūn)}{خمسون (khamsūn)}]{\ar{خمسون} (khamsūn) — fifty}

\label{lesson:AR-C97-khamsun}



\begin{quote}
Before the new one: say the Arabic for thirty, then the Arabic for forty.
\end{quote}

\subsection*{You'll want to know: \ar{خمسون}}
\textbf{\ar{خمسون}} — \emph{khamsūn} — \textquotedblleft{}fifty\textquotedblright{}.

\begin{cousinweb}
\textbf{\ar{خمسون}} (\emph{khamsūn}) is built on the native Arabic root \emph{kh-m-s}, \textquotedblleft{}five\textquotedblright{} (Lane). \emph{Khamsūn} is \emph{khams}, \textquotedblleft{}five\textquotedblright{}, with the tens ending \emph{-ūn}.
\end{cousinweb}

\begin{grammarlens}[title={What you've built}]
One more word for this chapter.
\end{grammarlens}

\subsection*{Your turn}
\begin{itemize}
\raggedright
  \item \emph{Say it:} \emph{khamsūn}
  \item \emph{Say it:} \emph{khamsūn}, once more
  \item \emph{Say it:} say \emph{khamsūn}
  \item \emph{Recall:} say the Arabic for thirty, then the Arabic for forty, then say \emph{khamsūn} again
\end{itemize}

\begin{tcolorbox}[breakable,colback=teal!4,colframe=teal!35!black,title={Before you move on}]
What does \ar{خمسون} mean? (\textquotedblleft{}Fifty\textquotedblright{}.) Say it once more.
\end{tcolorbox}
\section[\texorpdfstring{\ar{ألف} (alf)}{ألف (alf)}]{\ar{ألف} (alf) — a thousand}

\label{lesson:AR-C97-alf}



\begin{quote}
Before the new one: say the Arabic for forty, then the Arabic for fifty.
\end{quote}

\subsection*{You'll want to know: \ar{ألف}}
\textbf{\ar{ألف}} — \emph{alf} — \textquotedblleft{}a thousand\textquotedblright{}.

\begin{cousinweb}
\textbf{\ar{ألف}} (\emph{alf}) is built on the native Arabic root \emph{ʾ-l-f}, \textquotedblleft{}a thousand; to be familiar with\textquotedblright{} (Lane). Hebrew \emph{elef}, \textquotedblleft{}thousand\textquotedblright{}, is its cousin.
\end{cousinweb}

\begin{grammarlens}[title={What you've built}]
One more word for this chapter.
\end{grammarlens}

\subsection*{Your turn}
\begin{itemize}
\raggedright
  \item \emph{Say it:} \emph{alf}
  \item \emph{Say it:} \emph{alf}, once more
  \item \emph{Say it:} say \emph{alf}
  \item \emph{Recall:} say the Arabic for forty, then the Arabic for fifty, then say \emph{alf} again
\end{itemize}

\begin{tcolorbox}[breakable,colback=teal!4,colframe=teal!35!black,title={Before you move on}]
What does \ar{ألف} mean? (\textquotedblleft{}A thousand\textquotedblright{}.) Say it once more.
\end{tcolorbox}
